\chapter*{外文翻译}
\addcontentsline{toc}{chapter}{外文翻译}

\section*{原文}

Diffusion models represent a class of generative models that learn to reverse a gradual noising process. Given a data distribution $q(x_0)$, the forward process progressively adds Gaussian noise over $T$ timesteps, producing a sequence of increasingly noisy latents $x_1, x_2, \ldots, x_T$. The reverse process then learns to denoise, step by step, recovering the original data from pure noise.

The key insight is that if the noise added at each step is small enough, the reverse process can be approximated by a neural network trained to predict the noise component. This leads to stable training dynamics and high-quality sample generation, particularly for image synthesis tasks.

\section*{译文}

扩散模型是一类通过学习逆转逐步加噪过程的生成模型。给定数据分布 $q(x_0)$，前向过程在 $T$ 个时间步上逐步添加高斯噪声，产生一系列噪声逐渐增大的隐变量 $x_1, x_2, \ldots, x_T$。逆向过程则学习逐步去噪，从纯噪声中恢复原始数据。

其核心思想在于：如果每一步添加的噪声足够小，逆向过程可以用一个训练来预测噪声分量的神经网络来近似。这带来了稳定的训练动态和高质量的样本生成，尤其适用于图像合成任务。
